\begin{enumerate}[label=\textbf{1.1.\arabic*.},ref=1.1.\arabic*]
\setlength{\itemsep}{0.35em}
\item \label{formula-1-1-1} \textbf{Quadratic equations.} For $ax^2+bx+c=0$, $\Delta=b^2-4ac$ and $x=\frac{-b\pm\sqrt\Delta}{2a}$. The roots are rational when $\Delta$ is a non-negative rational square; $\alpha+\beta=-b/a$ and $\alpha\beta=c/a$.
\item \label{formula-1-1-2} \textbf{Section formula.} If $AP:PB=m:n$, then
\begin{equation*}\tag{1.1.2.1}
P=\left(\frac{nx_1+mx_2}{m+n},\frac{ny_1+my_2}{m+n}\right).
\end{equation*}
\item \label{formula-1-1-3} \textbf{Tangents to a circle.} $OQ\perp PQ$ at the point of contact, and tangents from an external point are equal. Hence $PQ=\sqrt{OP^2-r^2}$.
\item \label{formula-1-1-4} \textbf{Similarity.} For similar triangles, corresponding side ratios are equal; the ratio of their areas is the square of the ratio of corresponding sides.
\item \label{formula-1-1-5} \textbf{Step-deviation mean.} With $u_i=(x_i-a)/h$ and frequencies $f_i$,
\begin{equation*}\tag{1.1.5.1}
\bar{x}=a+h\frac{\sum f_i u_i}{\sum f_i}.
\end{equation*}
\item \label{formula-1-1-6} \textbf{Cone and hemisphere.}
\begin{equation*}\tag{1.1.6.1}
V_{\rm cone}=\frac13\pi r^2h,\qquad l=\sqrt{r^2+h^2},\qquad
S_{\rm cone,curved}=\pi rl,\qquad V_{\rm hemisphere}=\frac23\pi R^3.
\end{equation*}
\item \label{formula-1-1-7} \textbf{Arithmetic progressions.}
\begin{equation*}\tag{1.1.7.1}
a_n=a+(n-1)d,\qquad S_n=\frac n2\bigl(2a+(n-1)d\bigr),\qquad a_n=S_n-S_{n-1}.
\end{equation*}
\item \label{formula-1-1-8} \textbf{Probability.} $P(E)=\frac{n(E)}{n(S)}$ for equally likely outcomes; $P(E^c)=1-P(E)$. For independent events, $P(E\cap F)=P(E)P(F)$.
\item \label{formula-1-1-9} \textbf{Trigonometric identities.}
\begin{equation*}\tag{1.1.9.1}
\sin^2\theta+\cos^2\theta=1,\qquad 1+\tan^2\theta=\sec^2\theta,\qquad
\tan\theta=\frac{\sin\theta}{\cos\theta},\qquad\sin\theta=\frac{\text{opposite}}{\text{hypotenuse}},\quad\cos\theta=\frac{\text{adjacent}}{\text{hypotenuse}}.
\end{equation*}
\item \label{formula-1-1-10} \textbf{Sector of a circle.} If a circle is divided into $n$ equal sectors, one sector has area $\pi r^2/n$ and arc length $2\pi r/n$; its perimeter is $2r+2\pi r/n$.
\item \label{formula-1-1-11} \textbf{Coordinate geometry.} $AB=\sqrt{(x_2-x_1)^2+(y_2-y_1)^2}$; a circle centred at $(h,k)$ with radius $r$ has equation $(x-h)^2+(y-k)^2=r^2$.
\item \label{formula-1-1-12} \textbf{Parabola vertex.} For $y=ax^2+bx+c$ with $a\ne0$, the vertex has $x$-coordinate $-b/(2a)$ and $y$-coordinate $f(-b/(2a))$.

\item \label{formula-1-1-13} \textbf{Pairs of linear equations.} The lines $a_1x+b_1y+c_1=0$ and $a_2x+b_2y+c_2=0$ are inconsistent when $\frac{a_1}{a_2}=\frac{b_1}{b_2}\ne\frac{c_1}{c_2}$.
\item \label{formula-1-1-14} \textbf{HCF and LCM.} For positive integers $a,b$, $\operatorname{HCF}(a,b)\operatorname{LCM}(a,b)=ab$; the HCF can be found by prime factorization or Euclid's algorithm.
\end{enumerate}
